%قبل از شروع مسئله ترجیحاً چیزی ننویسید. اما اگر خواستید بنویسید، در قالب لیستی (فقرات و شمارش) بنویسید.
%برای اضافه کردن نمره، بعدی از شروع‌مسئله، [x] را اضافه کنید و به جای x  نمره مورد نظر را قرار دهید. برای زیرمسئله‌ها هم به همین صورت است.
%\شروع‌مسئله[10]
\شروع‌مسئله [12]
%\متن‌سیاه{سوال عملی اول:}\\
%\متن‌سیاه{دراپ‌تست}\\
محمد در یک مسابقه‌ی معلومات علمی شرکت کرده است. مسابقه شامل \متن‌لاتین{n} سوال است. محمد سوال \متن‌لاتین{i} ام را به احتمال \متن‌لاتین{p[i]} درصد درست پاسخ می‌دهد و به احتمال \متن‌لاتین{100 - p[i]} درصد غلط پاسخ می‌دهد. هم چنین اگر محمد سوال \متن‌لاتین{i} ام را درست پاسخ دهد، به اندازه‌ی \متن‌لاتین{c[i]} تومان جایزه می‌گیرد. حال روند مسابقه بدین شکل است که سوال ها به ترتیب از محمد پرسیده می شود. هر بار اگر محمد سوال را درست پاسخ دهد، جایزه آن سوال را می‌گیرد و در صورت وجود، به سوال بعد می‌رود. اما اگر یک سوال را غلط پاسخ دهد، مسابقه تمام می شود و محمد جایزه هایی که تا کنون جمع کرده است را به خانه می‌برد. حال ما می‌توانیم برای کمک به محمد، ترتیب سوال ها را تغییر دهیم. هدف این است که پس از تغییر ترتیب سوال ها، امیدریاضی جایزه‌ای که محمد می‌برد ماکسیمم شود. حال محمد از شما پرسیده است که به ازای همه‌ی ترتیب های ممکن، ماکسیمم امیدریاضی جایزه چقدر است ؟

\متن‌سیاه{ورودی:}\\
ابتدا در یک خط تعداد سوالات یعنی \متن‌لاتین{n} می‌آید.\\
سپس \متن‌لاتین{n} خط می‌آید که هر کدام شامل دو عدد  \متن‌لاتین{c[i]} و  \متن‌لاتین{p[i]} است.\\
\متن‌سیاه{خروجی:}\\
در تنها سطر خروجی، عدد خواسته شده را با دقت ۳ رقم اعشار چاپ نمایید.\\
\متن‌سیاه{محدودیت‌ها:}\\
$1 \leq n \leq 10^{5}$\\
$0 \leq p[i] \leq 100$\\
$0 \leq c[i] \leq 10^{4}$\\

\متن‌سیاه{ورودی و خروجی نمونه:}\\

تست اول
\begin{latin}
	\begin{tabularx}{\textwidth}{|X | X|}
		\hline
		\textbf{stdin} &
		\textbf{stdout}\\
		\hline
		4 & 2.434\\
		100 1&\\
		80 1 &\\
		66 1&\\
		20 1&\\
		\hline
	\end{tabularx}
\end{latin}

تست دوم
\begin{latin}
	\begin{tabularx}{\textwidth}{|X | X|}
		\hline
		\textbf{stdin} &
		\textbf{stdout}\\
		\hline
		3 & 2.326\\
		50 1&\\
		33 6&\\
		70 1&\\
		\hline
		
	\end{tabularx}
\end{latin}
%}
% صورت اصلی مسئله را در این‌جا بنویسید.
%اگر مسئله چند قسمت داشت به این صورت عمل کنید:
%\شروع{زیرمسائل}
%\زیرمسئله[0] % مسئلۀ ۱آ
%\زیرمسئله % مسئلۀ ۱ب
%\زیرمسئله % مسئلۀ ۱پ
%\زیرمسئله % مسئلۀ ۱ت
%\پایان{زیرمسائل}

\solution{
	برای حل این سوال،‌ فرض می‌کنیم که ترتیبی از سوال ها بیشترین امیدریاضی جایزه را داشته باشد. واضح است سوالاتی که حتما درست پاسخ می‌دهیم را در ابتدای ترتیب قرار دهیم. ترتیب بین آنها نیز اهمیتی ندارد. حال این سوالات را از مجموعه‌ی سوالات حذف می‌کنیم و از این جا به بعد بقیه سوال‌ها را در نظر می‌گیریم. فرض می‌کنیم سوال \متن‌لاتین{i}-ام در این ترتیب، حایزه‌ای برابر $c_i$ داشته باشد و احتمال پاسخ‌دهی درست به آن نیز $p_i$ باشد. حال دو سوال پشت سر هم \متن‌لاتین{i}-ام و \متن‌لاتین{i+1}-ام را در نظر بگیرید. ترتیب جدیدی از سوالات را در نظر بگیرید که در آن، جای این دو سوال عوض شده باشد. می‌دانیم امیدریاضی جایزه‌ی بدست آمده حداکثر برابر با حالت قبل است. با نوشتن امیدریاضی در دو حالت و لحاظ کردن کوچکتر مساوی بودن امیدریاضی حالت دوم بدست می‌آوریم که :
$$p_{i+1}*(1-p_i)*c_{i+1} \le p_i*(1-p_{i+1})*c_i$$
حال دو طرف را بر $(1-p_i) * (1-p_{i+1})$ تقسیم می‌کنیم. در این صورت خواهیم داشت:
$$\frac{p_{i+1}*c_{i+1}}{1-p_{i+1}} \le \frac{p_i*c_i}{1-p_i}$$
 بنابراین بدست آورده ایم که در حالت بهینه ترتیب سوال‌ها، سوال ها بر حسب 
$$\frac{p*c}{1-p}$$
 به صورت نزولی مرتب شده است. پس کافی است ما سوال ها را بر حسب رابطه بالا نزولی مرتب کنیم تا به حالت بهینه برسیم. اردر این راه نیز از اردر مرتب‌سازی خواهد بود که اگر از مرتب‌سازی ادغامی استقاده کنیم، از 
$O(nlogn)$
 خواهد بود. در زیر کد این سوال را که توسط آقای میکائیل قربانی زده شده است می‌بینید:

}




%%%%%%%%%%%%%%%%%%%%%%%%%%%%%%%%%%%%%%%%%%%%%%%%%%%%%
% برخی دستورات مفید برای تمرین‌ها. %
%%%%%%%%%%%%%%%%%%%%%%%%%%%%%%%%%%%%%%%%%%%%%%%%%%%%%

% فرمول‌های نمونه:
%$f(x) = \Theta(x)$
%$T(x, y) \leq \log(x) + 2^y + \binom{2n}{n}$
% {\tt code\_function}


% لیست موارد بدون شماره‌گذاری:
%\شروع{فقرات}
%    \فقره مورد اول 
%        \شروع{فقرات}
%            \فقره مورد اول داخلی 
%                \شروع{فقرات}
%                    \فقره اولین مورد خیلی داخلی 
%                    \فقره دومین مورد خیلی داخلی 
%                \پایان{فقرات}
%            \فقره مورد دوم داخلی 
%        \پایان{فقرات}
%    \فقره مورد دوم 
%\پایان{فقرات}

% لیست موارد شماره‌گذاری شده:
%\شروع{شمارش}
%    \فقره یک مورد 
%    \فقره یک مورد دیگر 
%    \فقره باز هم یک مورد دیگر
%\پایان{شمارش}

% برای نوشتن یک فرمول چند خطی که یک شماره بخوره. برای حالت‌های دیگر به سایت زیر مراجعه کنید:
% https://www.overleaf.com/learn/latex/Aligning_equations_with_amsmath
%\begin{align} 
%    \begin{split}
%        (x+y)^3 &= (x+y)^2(x+y) \\
%                &= (x^2+2xy+y^2)(x+y) \\
%                &= (x^3+2x^2y+xy^2) + (x^2y+2xy^2+y^3) \\
%                &= x^3+3x^2y+3xy^2+y^3
%    \end{split}                                 
%\end{align}

% نوشتن ماتریس و آرایه دوبعدی
%\begin{align}
%    A = 
%    \begin{bmatrix}
%        A_{11} & A_{21} \\
%        A_{21} & A_{22}
%    \end{bmatrix}
%\end{align}

% اضافه کردن تصویر با شرح. برای اضافه کردن تصویر بدون شرح، فقط دستور درج‌تصویر را استفاده کنید.
%\شروع{شکل}[h!]
%\تنظیم‌ازوسط
%  \درج‌تصویر[پهنا=0.5\پهنای‌متن]{img.jpg}
%\شرح{شرح تصویر}
%\برچسب{تصویر:۱}
%\پایان{شکل}
%\پایان{زیرمسائل}

%کد زیر را ببینید.
%\lstinputlisting{source_filename.py}
\begin{latin}
\begin{lstlisting}
array = []
n = int(input())
for i in range(0, n):
    possibility, bonus = [int(value) for value in input().split()]
    if possibility == 0 or bonus == 0:
        continue
    possibility /= 100
    array.append(((1 - possibility) / (possibility * bonus), possibility, bonus))

array.sort()
prefixPossibility = 1
answer = 0
for (sort_value, possibility, bonus) in array:
    prefixPossibility *= possibility
    answer += prefixPossibility * bonus

print('%.3f'% answer)

\end{lstlisting}
\end{latin}